\section{Additional results}\label{app:results}
\textbf{Paired comparison.}
Table~\ref{tab:wins} extends Section~\ref{sec:exp-c3} to every relaxation $\delta$ tried.
\begin{table}[h]
\caption{Paired comparison of \est{REG-L} with each competitor over 100 (response, $p$, $d$) combinations on the same draws. A win is a 95\% paired-bootstrap interval of the RMSE difference below zero, a loss above zero. ``Declared'' $\delta$ is the smallest on the grid with at least 50 expected exposed units per arm; ``oracle'' $\delta$ is the best in each cell and uses the truth.}
\label{tab:wins}
\centering
\input{../tables/tab_wins.tex}
\end{table}

\textbf{Observed-graph error.}
If a fraction $q$ of the edges the analyst uses are wrong (that many true edges dropped, as many random pairs added) while outcomes follow the true graph, the bias of \est{REG-L} in the one-hop world is $-q\gamma$ (normalised $-0.055,-0.13,-0.26,-0.38$ at $q=0.1,0.25,0.5,0.75$) and the standard deviation does not change; 88\% of the squared error at $q=0.25$ is bias. The error stays below \est{DM} at 0.50 at every $q$ tested, on ER, SBM, BA and WS alike (bias within 0.012 across graphs). The oracle form (\est{REG-O}) and \est{REG-Q} have the same error on the wrong graph; \est{REG-G}, which adds the two-hop exposure, is better than \est{REG-L} only when a two-hop term is present and $q\le0.25$ (0.22, 0.21, 0.27 against 0.25, 0.28, 0.32 at $q=0,0.1,0.25$) and worse in the one-hop world for $q\ge0.25$: no regression recovers an exposure that is not measured. A missed two-hop term and a wrong graph do not add (bias 0.245 at $q=0$ and 0.44 at $q=0.75$, not 0.62). The exact projection bias computed on the perturbed features matches the measured bias (correlation 0.9995). The perturbation is uniform edge rewiring; errors correlated with homophily or with the outcome, and missing edges only, are not tested.

\begin{figure}[h]
\centering
\includegraphics[width=\linewidth]{fig6a_misspec_map.pdf}
\caption{RMSE$/(|\beta|+|\gamma|)$ against the fraction $q$ of observed edges that are wrong, for one-hop and one-plus-two-hop interference (median over four graphs; band, median of the cell 95\% intervals; 100 replicates). \est{DM} is at 0.5.}
\label{fig:misspec}
\end{figure}
\textbf{Cluster scale and staggered rollout.}
Under graph cluster randomisation at $p=\tfrac12$, $n=2000$, $d=20$ (112 cells, 150 replicates) the Louvain resolution sets the cluster scale. The \est{REG-L} RMSE at the best resolution is 0.20--0.83 of its Bernoulli value on SBM ($\mu=0.8$) and WS graphs, where clusters hold at least 79\% of the edges, and 0.80--0.98 of it on ER graphs for the smooth responses; the threshold bias falls from 1.2 to 0.05 on the SBM with $\mu=0.8$. The U-shape in the cluster scale is mostly its right arm: coarsening past the community scale restores the Bernoulli error (threshold: 0.09, 0.93 and 1.11 at resolutions 5, 15 and 50 on the same SBM), and the left arm, too few clusters, appears only where the Louvain partition is coarse (Figure~\ref{fig:gcr}). The best resolution is taken in hindsight and differences below about 20\% between neighbouring resolutions are not resolved. Where the cluster scale is right, \est{DM} catches up with \est{REG-L} only on the threshold response.
Repeating $T+1$ independent Bernoulli measurements lowers the standard deviation (0.071 to 0.026 at $d=10$) and leaves the bias unchanged (threshold: $+0.83$ at every $T$). A staggered rollout, which observes the corner, gives an unbiased direct estimator (sd 0.015), whereas a regression pooled across its stages has a bias that grows with the number of stages (threshold, $d=10$: 0.04, 0.13, 0.19 for $T=2,4,8$) because interior stages dominate the fit (Table~\ref{tab:stag}).
\begin{figure}[h]
\centering
\includegraphics[width=\linewidth]{fig6b_gcr_ucurve.pdf}
\caption{\est{REG-L} RMSE$/(|\beta|+|\gamma|)$ against the Louvain resolution under graph cluster randomisation, four mechanisms on four graphs ($n=2000$, $d=20$, $p=\tfrac12$, 150 replicates); dotted: Bernoulli assignment.}
\label{fig:gcr}
\end{figure}

\begin{table}[h]
\caption{Staggered rollout against repeated Bernoulli measurements at matched numbers $T+1$ of measurements (ER, $n=2000$, 150 replicates): bias and RMSE with 95\% intervals of \est{REG-L} pooled over repeated Bernoulli experiments, pooled over the stages of a rollout, and of the direct rollout estimator.}
\label{tab:stag}
\centering
\input{../tables/app_tab_stag.tex}
\end{table}

\textbf{Full mechanism map.}
Tables~\ref{tab:cmedian} and~\ref{tab:cworst} give the median and the worst of the 60 cells of Table~\ref{tab:map} with 95\% intervals for every estimator.
\begin{table}[h]
\caption{Median over six graphs of RMSE$/(|\beta|+|\gamma|)$ with 95\% interval, by estimator and mechanism (setting of Table~\ref{tab:map}).}
\label{tab:cmedian}
\centering
\input{../tables/app_tab_c5_median.tex}
\end{table}
\begin{table}[h]
\caption{As Table~\ref{tab:cmedian}, worst graph.}
\label{tab:cworst}
\centering
\input{../tables/app_tab_c5_worst.tex}
\end{table}

\textbf{Interval coverage and width.}
\begin{table}[h]
\caption{Coverage of the true effect and mean width ($\gate=2$) of 95\% intervals by class of response: inside the curvature rule (linear, concave, convex), violating it (threshold, convex two-hop), and the count response; number of cells in parentheses, brackets 95\% bootstrap over cells. The last two rows inflate the iid and sandwich half-widths by one constant so that their mean width over all cells equals that of the bias-aware interval.}
\label{tab:intervals}
\centering
\input{../tables/tab2_coverage.tex}
\end{table}

\textbf{Intervals by regime.}
\begin{table}[h]
\caption{As Table~\ref{tab:intervals}, by regime: variance-dominated ($|\text{bias}|\le0.5\,\mathrm{sd}$), between, and bias-dominated ($|\text{bias}|>2\,\mathrm{sd}$) cells.}
\label{tab:regime}
\centering
\input{../tables/app_tab_cover_regime.tex}
\end{table}

\textbf{The extrapolation score.}
\begin{figure}[h]
\centering
\includegraphics{fig7_d1_scatter.pdf}
\caption{The extrapolation score D1 against $|\text{bias}|$ of \est{REG-L} in the 480 evaluation cells of Section~\ref{sec:diagnostic}.}
\label{fig:d1scatter}
\end{figure}

\section{Reproducibility details}\label{app:repro}
\textbf{Data.} Outcomes are simulated on every graph. The two real graphs supply topology only: the Stanford3 Facebook network (11,586 nodes, 568,309 edges; \citealp{traud2012social}, from the Network Repository, \citealp{rossi2015network}; CC BY-SA) and the ego-Facebook network (4,039 nodes, 88,234 edges; \citealp{mcauley2012learning}, from SNAP, which asks for citation). Both are reduced to the largest connected component; a thinning parameter keeps each edge with a fixed probability to vary the degree. The graphs are not redistributed; the released code downloads them and checks their SHA-256 digests.
\textbf{Replicates and seeds.} Every cell is seeded from a hash of its configuration; replicate $r$ uses the $r$-th child of that seed sequence, and grids that extend another replay its draws so that estimators are compared on identical replicates. Replicate counts are 100--500 per cell, stated with each result; intervals are Monte-Carlo intervals over replicates unless a caption says otherwise.
\textbf{Compute.} CPU only: shared containers with two cores. The full set of grids took about $3.7\times10^5$ CPU-seconds (about 100 CPU-hours); the largest are the method-by-mechanism grid with boosted trees and the graph network ($2\times10^5$ CPU-seconds) and the diagnostic grid ($6\times10^4$); the others take at most $4\times10^4$ each.
